\documentclass[11pt,a4paper]{article}
\usepackage[italian]{babel}
\usepackage[margin=2.5cm]{geometry}
\usepackage{parskip}
\usepackage{amsmath, amssymb, amsthm}
\usepackage{booktabs}
\usepackage{array}
\usepackage{xcolor}
\usepackage{needspace}
\usepackage{enumitem}
\usepackage{listings}
\usepackage{tikz}
\usepackage[most]{tcolorbox}
\usepackage[hidelinks]{hyperref}
\usetikzlibrary{decorations.pathreplacing, shapes.geometric, arrows.meta, positioning}

% Niente vedove/orfane: i paragrafi non lasciano righe isolate a fine/inizio pagina
\widowpenalty=10000
\clubpenalty=10000

% Elenchi più compatti
\setlist{itemsep=2pt, parsep=0pt, topsep=4pt, partopsep=0pt}

% --- Riquadri con bordo nero, senza numero (si spezzano tra due pagine) ---
% Uso: \begin{definizione} ... \end{definizione}
%      \begin{teorema}[Nome] ... \end{teorema}  (il nome è facoltativo)
\tcbset{nota/.style={enhanced jigsaw, breakable, colback=white,
  colframe=black, boxrule=0.5pt, sharp corners}}
\NewTColorBox{definizione}{}{nota,
  before upper={\textbf{Definizione.}\ }}
\NewTColorBox{teorema}{o}{nota,
  before upper={\textbf{Teorema\IfValueT{#1}{ (#1)}.}\ }}
\NewTColorBox{proposizione}{o}{nota,
  before upper={\textbf{Proposizione\IfValueT{#1}{ (#1)}.}\ }}
\NewTColorBox{proprieta}{o}{nota,
  before upper={\textbf{Proprietà\IfValueT{#1}{ (#1)}.}\ }}
\NewTColorBox{assioma}{o}{nota, boxrule=1pt,
  before upper={\textbf{Assioma\IfValueT{#1}{ (#1)}.}\ }}

% --- Ambienti semplici (senza riquadro) ---
% Il titolo sta in cima al blocco, anche se il contenuto inizia con un riquadro o
% con codice; e non resta mai da solo in fondo a una pagina.
% Uso: \begin{esempio}[Titolo facoltativo] ... \end{esempio}
\newcommand{\newrunin}[2]{%
  \NewDocumentEnvironment{#1}{o}{%
    \par\Needspace{4\baselineskip}\addvspace{\medskipamount}\noindent
    \textbf{#2}\IfValueT{##1}{ (##1)}\textbf{.}\ \ignorespaces}%
  {\par\addvspace{\medskipamount}}}
\newrunin{esempio}{Esempio}
\newrunin{esercizio}{Esercizio}
\newrunin{osservazione}{Osservazione}

% V e F nelle tavole di verità
\newcommand{\V}{\textsf{V}}
\newcommand{\F}{\textsf{F}}

% --- Codice C ---
% Uso: \begin{codice} ... \end{codice}      codice C numerato
%      \begin{comando} ... \end{comando}    comandi da terminale
%      \begin{risultato} ... \end{risultato} output di un programma
%      \lstinline|printf("ciao");|           codice dentro al testo
\lstdefinestyle{c}{
  language=C,
  basicstyle=\ttfamily\small,
  keywordstyle=\color{blue}\bfseries,
  commentstyle=\color{gray}\itshape,
  stringstyle=\color{red!70!black},
  numbers=left,
  numberstyle=\tiny\color{gray},
  numbersep=8pt,
  columns=flexible,
  keepspaces=true,
  breaklines=true,
  showstringspaces=false,
  tabsize=4,
  morekeywords={include, define},
  % lettere accentate nei commenti
  literate={à}{{\`a}}1 {è}{{\`e}}1 {é}{{\'e}}1 {ì}{{\`i}}1
           {ò}{{\`o}}1 {ù}{{\`u}}1 {È}{{\`E}}1 {À}{{\`A}}1
}
\lstset{style=c}
\newtcblisting{codice}{listing only, nota, left=6mm,
  listing options={style=c}}
\newtcblisting{comando}{listing only, nota,
  listing options={basicstyle=\ttfamily\small, numbers=none, language={},
    columns=flexible, keepspaces=true, breaklines=true}}
\newtcblisting{risultato}{listing only, enhanced jigsaw, breakable,
  colback=black!5, colframe=black, boxrule=0.5pt, sharp corners,
  listing options={basicstyle=\ttfamily\small, numbers=none, language={},
    columns=flexible, keepspaces=true, breaklines=true}}

% --- Diagrammi di flusso (simboli standard) ---
\tikzset{
  terminale/.style = {ellipse, draw, minimum width=2.5cm, minimum height=0.9cm},
  io/.style        = {trapezium, trapezium left angle=70, trapezium right angle=110,
                      draw, minimum width=2.5cm, minimum height=0.9cm},
  processo/.style  = {rectangle, draw, minimum width=2.5cm, minimum height=0.9cm},
  decisione/.style = {diamond, draw, aspect=2, minimum width=2.5cm},
  freccia/.style   = {thick, -{Stealth}}
}

\title{Ricorsione lineare: un esempio con radice doppia}
\author{Marco Cerbella}
\date{Lezione 7 -- 12 ottobre 2026}

\begin{document}
\maketitle

\section{Esempio con radici di molteplicità diversa}

Applichiamo il Teorema~B (radici con molteplicità) alla successione ricorsiva lineare omogenea a coefficienti costanti di grado $k = 3$
\[
\begin{cases}
a_0 = 2 \\
a_1 = -1 \\
a_2 = -8 \\
a_n = 5a_{n-1} - 3a_{n-2} - 9a_{n-3} & (n \geq 3).
\end{cases}
\]

\paragraph{Polinomio caratteristico.} L'equazione caratteristica è
\[
x^3 - 5x^2 + 3x + 9 = 0 .
\]

\paragraph{Radici.} Il polinomio è monico a coefficienti interi, quindi le radici razionali si cercano tra i divisori di $9$: $\pm 1, \pm 3, \pm 9$. Si trova che $3$ e $-1$ sono radici e
\[
(x - 3)^2 (x + 1) = 0 .
\]
\begin{itemize}
    \item $\lambda_1 = 3$ è radice con molteplicità algebrica doppia: $m_a(3) = 2$ ($= m_1$);
    \item $\lambda_2 = -1$ è radice semplice: $m_a(-1) = 1$ ($= m_2$);
    \item la somma delle molteplicità è $m_a(3) + m_a(-1) = 3 = k$, come deve essere.
\end{itemize}

\paragraph{Forma generale.} Per il Teorema~B
\[
a_n = (c_{1,0} + c_{1,1}\, n)\, 3^n + c_{2,0}\, (-1)^n .
\]

\paragraph{Condizioni iniziali.} Per trovare le costanti si impongono le condizioni iniziali (formula $=$ dato):
\[
\begin{cases}
n = 0: & c_{1,0} + c_{2,0} = 2 \\
n = 1: & (c_{1,0} + c_{1,1})\cdot 3 + c_{2,0}\,(-1) = -1 \\
n = 2: & (c_{1,0} + 2c_{1,1})\cdot 3^2 + c_{2,0}\,(-1)^2 = -8
\end{cases}
\]
Dalla prima equazione $c_{1,0} = 2 - c_{2,0}$. Sostituendo nelle altre due:
\[
\begin{cases}
6 - 3c_{2,0} + 3c_{1,1} - c_{2,0} = -1 \\
18 - 9c_{2,0} + 18c_{1,1} + c_{2,0} = -8
\end{cases}
\iff
\begin{cases}
-4c_{2,0} + 3c_{1,1} = -7 \\
-8c_{2,0} + 18c_{1,1} = -26
\end{cases}
\]
Dalla prima, $c_{1,1} = -\tfrac{7}{3} + \tfrac{4}{3} c_{2,0}$; sostituendo nella seconda
\[
-8c_{2,0} - 42 + 24c_{2,0} = -26 \iff 16 c_{2,0} = 16 \iff c_{2,0} = 1 .
\]
Quindi
\[
c_{2,0} = 1, \qquad c_{1,1} = -\tfrac{7}{3} + \tfrac{4}{3} = -1, \qquad c_{1,0} = 2 - 1 = 1 .
\]

\paragraph{Formula chiusa.}
\[
a_n = (1 - n)\, 3^n + (-1)^n .
\]

Verifica: $a_0 = 1 + 1 = 2$, $a_1 = 0 - 1 = -1$, $a_2 = (-1)\cdot 9 + 1 = -8$ e, dalla ricorrenza, $a_3 = 5(-8) - 3(-1) - 9\cdot 2 = -55 = (1-3)\cdot 27 - 1$.

\end{document}
